\cotefile{156-9}
\scanpage{126}
\mspage{120}
% archive sheet number: 125
% Lines 1--13; continuation from the preceding scan.
m'\deleted{\unsure{d}}\unsure{avait-il|aurait-il} trompé, pour une fois ? Il faudrait que j'y sois très attentif, chaque fois que je ferai appel à des considérations de support. Je voudrais p. ex. donner aux supports \unsure{attachés} à une figure, et leurs \unsure{relations} aux supports \unsure{attachés} à une subdivision, un ``sens'' qui ne soit affecté par le fait que les deux subdivisions ne sont ``\unsure{commensurables}'' \unsure{!!|''} --- ce ``souci'' un peu vague peut-il \deleted{\unsure{se}} préciser ?

% Lines 14--21.
On pourrait bien sûr songer à travailler plutôt avec le sous-ens. $\Sigma_{\unsure{\mathrm{ens}|\mathrm{en}}}$ de $\Sigma=\Sigma_A$, \unsure{formé} des supports \unsure{en.} constructibles, avec l'opération $\unsure{\mathcal{C}|\complement}$ (partout définie) et $\operatorname{Sup}$ et $\operatorname{Inf}$ (partiellement \unsure{définies|définie} seulement). \unsure{Cet|Peut} ens. \unsure{admet-il} une construction \unsure{indépendante} de la considé%
\scanpage{127}%
\mspage{121}%
% archive sheet number: 126
% Lines 1--16; the opening completes considération on scan 126.
ration préalable de $\Sigma$, le mal-aimé ? Par exemple, les supports \inserted{\unsure{constr.}} \uline{fermés} (\deleted{\unsure{associés}} \inserted{\unsure{relatifs}} à un magasin-\unsure{étalons|ateliers} $\unsure{\mathcal{F}|F}$) \unsure{peuvent-ils|pourront-ils} s'obtenir, en prenant l'ensemble ordonné déduit de $(\mathcal{F},\ll)$ en ``inversant'' les ``flèches'' $\unsure{\ll|\longleftarrow}$ (sur les subdivisions) ? Je ne suis pas $=$ sûr que la cat. des \unsure{fractions} obtenue soit bien un ens. ordonné, et encore moins (\unsure{même} s'il l'était) qu'il \deleted{\unsure{coïncide avec}} s'identifie à un sous-ens. ordonné de $\Sigma$. (Ce n'est pas là nécessairement un mal, si $\Sigma$ devait être \unsure{largué|légué} aux profits et pertes\ldots).

% Lines 17--23.
Je dirais quand-même, pour me \unsure{rassurer}, que dans les cas où $\Sigma$ est booléen (ateliers \unsure{ponctuaires}, et peut-être $=$ les ateliers fidèles ?, ou tout au moins \unsure{ainsi|aussi} les \unsure{maquettes}), il n'y a rien à lui reprocher. Ça inclut quand même déjà tous les cas qu'on peut
\scanpage{128}
\mspage{122}
% archive sheet number: 127
% Lines 1--7; continuation from scan 127.
considérer \unsure{tant} soit peu ``classiques''. Ce sont mes ``\unsure{cordes}'', \unsure{pseudo-simplexes} topologiques et tutti quanti, associés à l'intuition des ``lieux'' remplaçant les ``points'', qui causent tous les ennuis pour l'instant, et les seuls ennuis.

% Lines 8--16. A vertical emphasis rule stands to the left of this paragraph.
\noindent\begin{tabular}{|p{0.9\linewidth}}
Mais même dans le cas des ateliers \unsure{ponctuaires}, il serait intéressant de voir si $\Sigma_{\unsure{\mathrm{ens}|\mathrm{en}}}$ n'admet pas une description combinatoire directe, dans un esprit plus proche d'une ``géométrie des formes'' --- alors que la \unsure{Dfation|Dftion} a des forts relents de topologie générale !
\end{tabular}
\par\centerline{\rule{2cm}{0.4pt}}

Prochaines choses à regarder
% Right margin beside the end of the preceding paragraph and the list.
\margin{pour ``boucler'' l'aspect purement ensembliste --- il serait temps !}
% margin note written slanted, to the left of the list; the word morphismes extends above item a).
\margin{et avant encore, peut-être retourner sur les \uline{morphismes} d'ateliers, et leur \unsure{relation} aux supports}
\begin{enumerate}[label=\alph*)]
\item Subdivisions --- enfin !
\item Ateliers modérés
\item Axiomes de saturation. Saturation d'un atelier, ou d'un atelier-magasin
\end{enumerate}
\centerline{\ldots}

\scanpage{130}
\mspage{1}
% archive sheet number: 129
12.7.86 \hfill CF IX bis

% margin note written slanted, rising along the upper left edge.
\margin{Comparer CF $\overline{\underline{\unsure{\mathrm{XIII}|\mathrm{XII}}}}$ p. 46 ff. et aussi p. 27 ff.}
\uline{Morphismes d'ateliers}

Je suis intéressé par trois types de situations, que je voudrais pouvoir exprimer par une notion de morphismes

1\textsuperscript{o}) Cas d'un atelier \inserted{$A$} et d'un sous-atelier $A'$, i.e. (cf CF $\overline{\underline{\mathrm{IX}}}$ p. 27) une partie $A'$ satisfaisant la condition
\[
\begin{aligned}
&\text{\uline{Sous at}}\quad \text{Si }X\overset{\circ}{\ll}Y'\triangleleft Y,\quad\text{avec }X,Y\in A',\ \text{alors}\\
&Y'\in A'.\quad\unsure{!}
\end{aligned}
\]

Cela implique que, muni des relations induites par $A$ $\unlhd$, $\overset{\circ}{\ll}$, $|\unsure{\circ|\bullet}|$, $A'$ est un atelier, et que de plus, la relation $\ll_{A'}$ est induite par $\ll_A$ (ce qui est \unsure{assuré} précisément par Sous at). Je rappelle les exemples principaux

\uline{Ex 1} Soit \unsure{\textsuperscript{*}} $\unsure{\underline{\varphi}|\underline{\Phi}}\inserted{\subset A}$ une préfigure de $A$, i.e. une partie telle que $\overline{\underline{\varphi}}^{(\unlhd)}$ soit une figure, i.e. telle que $X,\deleted{\unsure{X}}\unsure{Y}\in\underline{\varphi}\Rightarrow X\unsure{\asymp|\bowtie}Y$. Alors $\operatorname{Omb}^{\circ}(\underline{\varphi})=\{X\in A\mid X\overset{\circ}{\ll}\underline{\varphi}\}$ est un sous-atelier. (loc. cit)

\uline{Ex} \uline{1 bis} $\underline{\varphi}$ une \deleted{\unsure{pré}}figure de $A$, alors $\underline{\varphi}$ lui-même\inserted{!} est un \unsure{sous-atelier}, (qui est une \uline{maquette} pour la structure induite). (En effet, si $X\overset{\circ}{\ll}Y'\unlhd Y$, avec $X,Y\in\underline{\varphi}\inserted{\subset\overline{\underline{\varphi}}}$, on a $Y'\in\overline{\underline{\varphi}}$, \unsure{mais aussi} $X=Y'$, d'où $Y'\in\underline{\varphi}$.)

\scanpage{131}
\mspage{2}
% archive sheet number: 130
\uline{Ex 2} Soit $S\subset A$, considérons ``l'atelier résiduel strict'' formé des $X\in A$ tels que $X\mathrel{\Vert}S$, il est \deleted{\unsure{stable}} \inserted{fermé} pour $\ll$ et \unsure{saturé} pour $\unlhd$, donc c'est un sous-atelier.

\uline{Exemple 2 bis} Considérons par contre le sous-ensemble $\operatorname{comp}^{\circ}(S)=\{X\in A\mid X\mathrel{|\unsure{\circ|\bullet}|}S\}$. Ce n'est pas en général un sous-atelier (loc. cit p. 28). Mais dans le cas où $S=\unsure{\underline{\varphi}|\underline{\Phi}}$ est une préfigure, considérons l'ensemble \deleted{\unsure{des}}
\[
\{X\in A\mid X\mathrel{|\bullet|}\underline{\varphi}\ \text{et}\ X\unsure{\asymp|\bowtie}\underline{\varphi}\}
\]
(atelier résiduel au sens large). C'est un sous-atelier. En effet, soit
\[
X\overset{\circ}{\ll}Y'\unlhd Y\qquad\text{avec }X,Y\in A',
\]
prouvons $Y'\in A'$, i.e.
\[
\forall Z\in\underline{\varphi},\qquad Y'\deleted{\unsure{\triangleright|\asymp}}\unsure{\asymp|\triangleright}Z,\quad\text{et}\quad Y'\mathrel{|\circ|}Z\ \text{i.e. }Y'\ne Z
\]
\unsure{La} relation \shortstack[l]{\deleted{$Y'\asymp Z$ \unsure{claire car} $Y\asymp Z$}\\[-0.4ex]\makebox[0.48\linewidth]{-\hfill-\hfill-\hfill-\hfill-\hfill-\hfill-\hfill-\hfill-}}, \inserted{$Y'\unlhd Y$,} si on avait $Y'=Z$, on aurait $X\overset{\circ}{\ll}Z$, or $X\in A'$ implique $X\mathrel{|\circ|}Z$, donc on aurait $X\mathrel{|\circ|}X$, absurde.
% Nine dash segments below the cancelled phrase are retained in the shortstack; their restoration status is uncertain.

\margin{$\Vert$}
Dans ces quatre cas, il faudrait vérifier si $\Sigma_{A'}\hookrightarrow\Sigma_A$ ?

\scanpage{132}
\mspage{3}
% archive sheet number: 131
2\textsuperscript{o}) Cas d'une application ensembliste \deleted{de} $f:E'\unsure{\longrightarrow|{-}}E$, ou d'espaces modérés\inserted{\ldots}, et de l'application
\[
\unsure{f^*}:\operatorname{Figures}(E)\longrightarrow\operatorname{Figures}(E')
\]
% margin note written slanted, rising along the left side of the display.
\margin{dans le cas d'une application modérée d'espaces modérés, \unsure{en prenant} les ``figures \unsure{modérées}''}
correspondante. Si $f$ n'est pas surjective, cette application \unsure{n'applique} pas $\unsure{\mathrm{Figel}\,\mathrm{ens}|\mathrm{Figelens}}\inserted{(E)=A}$ dans $\mathrm{Figel}\,\mathrm{ens}\inserted{(E')=A'}$, mais dans \deleted{$\mathrm{Figel}\,\mathrm{ens}\cup\{\unsure{o|\varnothing|0}$} $A'\cup\{\varnothing\}$.
\inserted{$\leftarrow$ figure vide}
% The insertion's leftward arrow points to the empty-set symbol.

Cette application est compatible avec $\unlhd$, $\overset{\circ}{\ll}$, $|\unsure{\circ|\bullet}|$, mais en général elle n'est pas injective\ldots\ . Elle a tendance à passer aux supports, et alors à commuter à tout ($\operatorname{Sup}$, $\operatorname{Inf}$, $C$).

3\textsuperscript{o}) Cas d'une application ensembliste $g:E\longrightarrow E'$\unsure{,|*} injective, et application
\[
g_*:\deleted{\mathrm{Figel}}\ A\longrightarrow A',\qquad\text{également injective.}
\]
Commute encore à : $\unlhd$, $\overset{\circ}{\ll}$, $|\circ|$, mais en passant aux supports, commute à $\underset{i\in I}{\operatorname{Sup}}$ ($I$ quelc\ldots), $\underset{i\in I}{\operatorname{Inf}}$ ($I\ne\varnothing$), $S\cap T$, mais pas à $C$ si $g$ non bijective.

\scanpage{133}
\mspage{4}
% archive sheet number: 132
Cas commun à 2\textsuperscript{o} et 3\textsuperscript{o} :

correspondance
\diagram{\begin{tikzcd}[ampersand replacement=\&,column sep=large,row sep=large]
\Gamma \arrow[r,hook,"\unsure{f'|g'}"] \arrow[d,"f"'] \& E' \\
E \&
\end{tikzcd},}
telle que $\Gamma\unsure{\longrightarrow|{-}}E'$ soit \deleted{\unsure{c}}injectif\inserted{!!}. Dans ce cas qui \unsure{généralise} 2\textsuperscript{o} et 3\textsuperscript{o} (en fait, une \unsure{composition} des deux) on a encore
\[
A\longrightarrow A'\cup\{\varnothing\}.
\]

% margin note written slanted, rising along the left edge and overlapping the main text.
\margin{Le test si on est dans le cas 2\textsuperscript{o}, est $\deleted{\unsure{f_*|\alpha_*}}\unsure{\alpha_*|f_*}(1)=1$. Le test pour le cas 3\textsuperscript{o} est quoi ?? Pas l'injectivité sur les supports (C'est le fait que $f$ soit injectif, \unsure{sans} être néc. iso) ni la \unsure{fidélité de} la conservation de $|\unsure{\circ|\bullet}|$ (\unsure{$=$} \unsure{univoque}), $\unsure{C|\mathcal{C}}$ !! alors quoi ? Le fait que $\beta_*$ \unsure{commute à} $\alpha_*$. (En plus de l'injectivité de l'application sur les supports et de la conservation de $|\circ|$, ou de $S\cap T$, \unsure{$S\cap C T$|$S\cap\complement T$}\ldots\unsure{$c$|$C$}).}

Il ne semble pas y avoir d'application directe de figures, pour une application $g:E\longrightarrow E'$ qui ne soit injective.

Dans le cas d'un morphisme d'ateliers quelconque, il faut voir si se \unsure{construit|induit} tjs une ``correspondance'' sur les \unsure{spatios} associés --- d'où un $\unsure{\alpha_*|\alpha^*}$, $\unsure{\beta_*|\beta^*}$.

